\chapter{Introdução}
\label{cap:introducao}


A classificação atribui um rótulo a cada instância de um conjunto de dados. Em muitas
aplicações, porém, o interesse recai sobre a proporção de cada classe em um conjunto de
instâncias não rotuladas, e a identificação individual é apenas um meio de chegar a
ela \cite{forman2005}. A tarefa de estimar essa proporção recebe o nome de
quantificação \cite{esuli2023}.

A diferença entre as duas tarefas tem consequência prática. Contar os rótulos atribuídos
por um classificador acumula os erros dele, e o resultado carrega um viés que cresce
conforme a proporção real se afasta daquela observada no treinamento. Por isso foram
propostos métodos próprios para a quantificação, com correções que vão da compensação
das taxas de acerto do classificador ao casamento entre distribuições de escores.

A tarefa aparece em áreas distintas. No monitoramento de insetos vetores de doenças,
interessa a proporção de cada espécie capturada em uma região, informação que orienta
ações de controle sem exigir a identificação de cada inseto \cite{maletzke2019}. Em
redes sociais, interessa a fração de menções positivas e negativas sobre um tema. Em
ambos os casos, a estimativa da proporção é o produto final.

A literatura da área acumula dezenas de métodos, e as implementações disponíveis em
Python estão distribuídas em bibliotecas independentes. Cada uma adota convenções
próprias para os dados, para a construção dos quantificadores e para a avaliação, de
modo que reproduzir um experimento em outra biblioteca exige reescrever o código que o
executa. Nenhuma delas cobre o conjunto de métodos propostos na literatura, e a escolha
da biblioteca passa a restringir os métodos ao alcance de quem a utiliza.

Esse cenário motiva duas perguntas. A primeira trata da possibilidade de reunir os
métodos sob uma interface única, com a mesma convenção de uso da biblioteca de
aprendizado de máquina mais difundida em Python. A segunda trata da garantia de que uma
implementação nova produza o mesmo resultado das existentes, uma vez que a adoção de uma
biblioteca depende da confiança em suas estimativas.

Neste trabalho, é proposta a \texttt{mlquantify}, uma biblioteca de quantificação
construída sobre a interface do \textit{scikit-learn}, com os métodos, os protocolos de
avaliação e as métricas de erro da área. A biblioteca é comparada experimentalmente com
a \verb'QuaPy', a \verb'QuantificationLib' e a \verb'qunfold', nos métodos que cada uma
implementa, quanto ao erro de estimativa e ao custo computacional.


\section{Objetivos}
\label{sec:objetivos}

O objetivo geral deste trabalho é desenvolver uma biblioteca de quantificação em Python
que reúna os métodos da literatura sob uma interface única e avaliá-la frente às
bibliotecas existentes.

Os objetivos específicos são os seguintes.

\begin{enumerate}
    \item Definir uma arquitetura capaz de acomodar métodos de exigências distintas sob
    uma mesma interface.
    \item Implementar os métodos de quantificação, os protocolos de avaliação e as
    métricas de erro descritos na literatura.
    \item Reunir e selecionar uma coleção de conjuntos de dados tabulares para a
    avaliação experimental.
    \item Verificar se o erro da biblioteca proposta é equivalente ao das bibliotecas de
    referência em um mesmo método.
    \item Medir o custo em tempo e em memória da biblioteca proposta em relação às
    demais.
\end{enumerate}


\section{Organização do Texto}
\label{sec:organizacao_texto}

O restante deste texto está organizado em cinco capítulos. No
Capítulo~\ref{cap:aprendizado_de_maquina}, são apresentados os conceitos de aprendizado de máquina
que sustentam a quantificação, com o aprendizado indutivo, as tarefas supervisionadas e
os métodos de avaliação. No Capítulo~\ref{chap:quantificacao}, a quantificação é
definida, junto com a taxonomia dos métodos, os métodos em si e os protocolos e métricas
usados para avaliá-los. No Capítulo~\ref{chap:materiais_metodos}, são descritas as
bibliotecas de referência e o método adotado no desenvolvimento e na avaliação da
biblioteca proposta. No Capítulo~\ref{cap:resultados}, são apresentadas a biblioteca
desenvolvida e os resultados da comparação. No Capítulo~\ref{cap:conclusao}, são
reunidas as conclusões, as limitações e as possibilidades de continuidade.
